%% ════════════════════════════════════════════════════════════════════════════
%% PLANTILLA DE INFORME INACAP
%% Compilar con: pdflatex -shell-escape plantilla_informe_inacap.tex
%% (se requiere -shell-escape por el paquete minted; ejecutar 2 veces para el índice)
%% Archivos necesarios en la misma carpeta: logo_inacap.png
%% ════════════════════════════════════════════════════════════════════════════
\documentclass[11pt,a4paper]{article}
\usepackage[utf8]{inputenc}
\usepackage[spanish]{babel}
\usepackage[T1]{fontenc}
\usepackage{geometry}
\usepackage{amsmath}
\usepackage{amsfonts}
\usepackage{amssymb}
\usepackage{booktabs}
\usepackage{hyperref}
\usepackage{enumitem}
\usepackage{listings}
\usepackage{titlesec}
\usepackage{setspace}
\usepackage{tikz}
\usepackage{xcolor}
\usepackage{eso-pic}
\usepackage{graphicx}
\usepackage{fancyhdr}
\usepackage{float}
\usepackage{minted}
\setlength{\headheight}{42pt} % Altura suficiente para 3 líneas

% ── Colores corporativos INACAP ─────────────────────────────────────────────
\definecolor{inacapRed}{RGB}{204,0,0}
\definecolor{lightGray}{RGB}{200,200,200}

% ── Variables del documento (editar solo aquí) ──────────────────────────────
\newcommand{\docTitulo}{Título del Informe}
\newcommand{\docEvaluacion}{Nombre de la Evaluación}
\newcommand{\docCarrera}{Nombre de la Carrera}
\newcommand{\docAsignatura}{Nombre de la Asignatura}
\newcommand{\docSeccion}{Código / Sección}
\newcommand{\docDocente}{Nombre del Docente}
\newcommand{\docEstudiante}{Nombre del Estudiante (o integrantes del grupo)}
\newcommand{\docFecha}{DD de Mes de AAAA}

% ── Cabecera de página ───────────────────────────────────────────────────────
\fancypagestyle{mainpage}{%
  \fancyhf{}
  % Línea roja bajo el encabezado
  \renewcommand{\headrule}{\color{inacapRed}\hrule width\headwidth height 1.5pt \vss}
  \renewcommand{\headrulewidth}{1.5pt}
  % Logo a la izquierda
  \fancyhead[L]{\includegraphics[height=1.1cm,keepaspectratio]{logo_inacap.png}}
  % Información a la derecha
  \fancyhead[R]{%
    \footnotesize\color{gray}%
    \docCarrera\\%
    \docAsignatura\\%
    \docSeccion}
  % Número de página centrado al pie
  \fancyfoot[C]{\thepage}
}

% ── Configuración de márgenes (contenido principal) ─────────────────────────
\geometry{top=2.5cm, bottom=2.5cm, left=3cm, right=3cm}

% ── Estilo de títulos ────────────────────────────────────────────────────────
\titleformat{\section}{\large\bfseries}{\thesection}{1em}{}
\titleformat{\subsection}{\normalsize\bfseries}{\thesubsection}{1em}{}

\title{\docTitulo}
\author{}
\date{}

% ════════════════════════════════════════════════════════════════════════════
\begin{document}

% ════════════════════════════════════════════════════════════════════════════
% PORTADA
% ════════════════════════════════════════════════════════════════════════════
\pagestyle{empty}
\newgeometry{top=0cm, bottom=2cm, left=2cm, right=2cm}

\AddToShipoutPictureBG*{%
\begin{tikzpicture}[remember picture, overlay]

  % --- Fondo gris claro franja superior ---
  \fill[lightGray!50]
    (current page.north west)
    rectangle
    ([yshift=-2.2cm] current page.north east);

  % --- Logo INACAP ---
  \node[anchor=north west, inner sep=0pt]
    at ([xshift=1.5cm, yshift=0cm] current page.north west)
    {\includegraphics[width=4cm, height=2.2cm, keepaspectratio]{logo_inacap.png}};

  % --- Rayas rojas decorativas derecha (10 rayas) ---
  \foreach \i in {0,1,2,3,4,5,6,7,8,9}{
    \pgfmathsetmacro{\xstart}{-3.5 + \i*0.28}
    \pgfmathsetmacro{\xend}{\xstart + 0.16}
    \fill[inacapRed]
      ([xshift=\xstart cm, yshift=-0.15cm] current page.north east)
      rectangle
      ([xshift=\xend cm, yshift=-2.05cm] current page.north east);
  }

  % --- Bloque rojo sólido esquina superior derecha ---
  \fill[inacapRed]
    ([xshift=-0.9cm, yshift=0cm] current page.north east)
    rectangle
    ([xshift=0cm, yshift=-2.2cm] current page.north east);

  % --- Líneas decorativas diagonales abajo-derecha ---
  \draw[lightGray, line width=0.8pt]
    ([xshift=-2cm, yshift=4.5cm] current page.south east)
    --
    ([xshift=1cm, yshift=3.0cm] current page.south east);
  \draw[lightGray!70, line width=0.5pt]
    ([xshift=-4cm, yshift=5.0cm] current page.south east)
    --
    ([xshift=-1cm, yshift=3.5cm] current page.south east);

\end{tikzpicture}%
}

% ── Espacio para la barra de encabezado ─────────────────────────────────────
\vspace{2.8cm}

\begin{center}
  \vspace{4cm}

  {\large\bfseries \docEvaluacion}

  \vspace{0.6cm}
\end{center}

\vspace*{\fill}

\begin{center}
    {\huge\bfseries \docTitulo}
\end{center}

\vfill
\noindent
\begin{tabular}{@{}p{5.5cm} p{9cm}@{}}
  \textbf{Asignatura:}              & \docAsignatura \\[0.4em]
  \textbf{Sección:}                 & \docSeccion \\[0.4em]
  \textbf{Nombre del docente:}      & \docDocente \\[0.4em]
  \textbf{Nombre de los}            & \\
  \textbf{integrantes del grupo:}   &
    {\color{inacapRed} \docEstudiante}   \\[0.15em]
\end{tabular}

\vspace{1.5cm}

\begin{center}
  \textbf{\docFecha}
\end{center}

% ════════════════════════════════════════════════════════════════════════════
% CONTENIDO DEL INFORME
% ════════════════════════════════════════════════════════════════════════════
\clearpage
\restoregeometry
\pagestyle{mainpage}

\vspace{1cm}

\tableofcontents
\newpage

% ── Introducción ─────────────────────────────────────────────────────────────
\section*{Introducción}
Texto de la introducción.

\subsection*{Subtítulo sin numerar}
\begin{enumerate}[label=\alph*)]
    \item Primer elemento.
    \item Segundo elemento.
\end{enumerate}

% ── Sección numerada ─────────────────────────────────────────────────────────
\newpage
\section{Título de la sección}

\subsection{Título de la subsección}
Texto de la subsección.

\begin{enumerate}[label=\alph*)]
    \item Elemento con lista literal.
    \item Otro elemento.
\end{enumerate}

\subsubsection*{Subsubsección sin numerar}
Texto.

% ── Tabla ────────────────────────────────────────────────────────────────────
\newpage
\subsection{Tabla comparativa}
\begin{table}[H]
\centering
\begin{tabular}{@{}p{3.5cm}p{5cm}p{5cm}@{}}
\toprule
\textbf{Criterio} & \textbf{Columna A} & \textbf{Columna B} \\ \midrule
Criterio 1 & \textbf{Etiqueta.} Descripción. & \textbf{Etiqueta.} Descripción. \\
Criterio 2 & \textbf{Etiqueta.} Descripción. & \textbf{Etiqueta.} Descripción. \\
Criterio 3 & \textbf{Etiqueta.} Descripción. & \textbf{Etiqueta.} Descripción. \\ \bottomrule
\end{tabular}
\end{table}

% ── Subsección con descripción/aplicación ────────────────────────────────────
\section{Otra sección}

\subsection{Elemento con descripción}
\textbf{Descripción:} Texto descriptivo.

\textbf{Aplicación en el caso:} Texto de aplicación.

% ── Bloques de código (minted) ───────────────────────────────────────────────
\newpage
\section{Ejemplos de código}

\subsection{Ejemplo 1}

\subsubsection*{Prompt / Entrada:}
\begin{minted}[breaklines]{python}
def funcion(param):
    return param
\end{minted}

\subsubsection*{Salida esperada:}
\begin{minted}[breaklines]{json}
{
  "clave": "valor"
}
\end{minted}

\begin{minted}[breaklines]{text}
Texto plano de salida.
\end{minted}

% ── Conclusiones ─────────────────────────────────────────────────────────────
\newpage
\section*{Conclusiones}
Párrafo de conclusión 1.

Párrafo de conclusión 2.

% ── Bibliografía ─────────────────────────────────────────────────────────────
\newpage
\section{Bibliografía}
\begin{enumerate}
\item \textbf{Autor (año).} \textit{Título de la obra}. Editorial. \url{https://ejemplo.com}
\item \textbf{Autor (año).} \textit{Título de la obra}. Editorial.
\end{enumerate}

\end{document}
